% titel: Blatt 4, A7: Welche Serie? (Wasserstoffspektrum, Wellenzahlen)
% skript: Mortimer Kap. 6 (Serien des Wasserstoffs)

% Teil I: Werkzeuge, kein Lösungsweg.

Die Linien des Wasserstoffatoms ordnen sich in Serien, benannt nach dem
unteren Niveau $n_f$ (Lyman $n_f=1$, Balmer $n_f=2$, Paschen $n_f=3$,
Brackett $n_f=4$, Pfund $n_f=5$). Zwei aufeinanderfolgende Linien gehören
zur selben Serie, wenn sie denselben Endzustand $n_f$ teilen.

\begin{werkzeugkasten}[Serienidentifikation über die Wellenzahl]
Wellenzahl einer Linie von $n_i \to n_f$:
\begin{equation}
  \tilde\nu = \frac{1}{\lambda} = R_H\left(\frac{1}{n_f^2}
  - \frac{1}{n_i^2}\right),
  \qquad R_H \approx \qty{1.097e7}{m^{-1}}.
\end{equation}
Zwei aufeinanderfolgende Linien derselben Serie haben dasselbe $n_f$ und
$n_i$ bzw.\ $n_i+1$. Für Balmer ($n_f=2$):
\begin{equation}
  \tilde\nu = R_H\!\left(\frac14 - \frac{1}{n_i^2}\right).
\end{equation}
Man kann durch Probieren der $n_f$ prüfen, welche Serie beide
gegebenen Werte reproduziert.

\textit{Skript:} Mortimer Kap. 6 (Serien, Rydberg-Formel), Übung 4/7.
\end{werkzeugkasten}

\begin{achtung}[Darauf achten]
\begin{itemize}
  \item Beide Wellenzahlen müssen zur \emph{gleichen} Serie gehören, d.\,h.
        gleichen $n_f$ und $n_i$-Werte haben, die sich um $1$ unterscheiden.
  \item Die Serien sind nach dem unteren Niveau benannt; nicht nach dem
        oberen.
  \item Ohne $R_H$ geht der Kürzest-weg-Test: Der Quotient bzw.\ die
        Differenz der Wellenzahlen legt die Serienzugehörigkeit fest.
  \item Werte in m$^{-1}$ einsetzen ($10^6$-Faktor beachten).
\end{itemize}
\end{achtung}
